\chapter{Model Evaluation and Selection}

\section{Training, Validation, and Testing}

In previous chapters, we developed models and optimization methods for learning from data. However, a crucial practical question remains:

\medskip

\noindent
\textbf{How do we evaluate and compare models in a reliable way?}

\medskip

A central challenge is that the true risk
\[
R(f) = \mathbb{E}_{(x,y)\sim P}[\ell(f(x), y)]
\]
is unknown, since the data-generating distribution $P$ is not accessible.

\medskip

\noindent
\textbf{Data splitting.}  
To approximate generalization performance, we divide the available dataset into three parts:

\begin{itemize}
    \item \textbf{Training set:} used to fit model parameters
    \item \textbf{Validation set:} used to tune hyperparameters
    \item \textbf{Test set:} used to estimate final performance
\end{itemize}

\medskip

\noindent
\textbf{Workflow.}
\begin{enumerate}
    \item Train model on training set
    \item Select hyperparameters using validation set
    \item Evaluate final model on test set
\end{enumerate}

\medskip

\noindent
\textbf{Key principle.}  
The test set must not be used during model selection, as this would lead to overly optimistic estimates.

\medskip

\noindent
\textbf{Overfitting to validation data.}  
Repeated tuning can indirectly overfit the validation set. This motivates more robust evaluation strategies, such as cross-validation.

\section{Cross-Validation}

Cross-validation provides a systematic way to estimate model performance using limited data.

\medskip

\noindent
\textbf{$k$-fold cross-validation.}
\begin{itemize}
    \item Partition the dataset into $k$ subsets (folds)
    \item For each fold:
    \begin{itemize}
        \item Train on $k-1$ folds
        \item Evaluate on the remaining fold
    \end{itemize}
    \item Average the results across folds
\end{itemize}

\medskip

\noindent
\textbf{Cross-validation estimate.}
\[
\hat{R}_{\text{CV}} = \frac{1}{k} \sum_{j=1}^k R_j,
\]
where $R_j$ is the validation error on fold $j$.

\medskip

\noindent
\textbf{Advantages.}
\begin{itemize}
    \item Efficient use of data
    \item More stable performance estimates
\end{itemize}

\medskip

\noindent
\textbf{Choice of $k$.}
\begin{itemize}
    \item Small $k$ (e.g., $5$): faster, higher bias
    \item Large $k$ (e.g., leave-one-out): lower bias, higher variance and cost
\end{itemize}

\medskip

\noindent
\textbf{Practical use.}  
Cross-validation is widely used for hyperparameter tuning, especially when datasets are small or moderate in size.

\section{Performance Metrics}

The choice of evaluation metric depends on the learning task.

\subsection{Regression Metrics}

\begin{itemize}
    \item \textbf{Mean Squared Error (MSE):}
    \[
    \frac{1}{n} \sum_{i=1}^n (f(x_i) - y_i)^2
    \]

    \item \textbf{Mean Absolute Error (MAE):}
    \[
    \frac{1}{n} \sum_{i=1}^n |f(x_i) - y_i|
    \]
\end{itemize}

\medskip

\noindent
\textbf{Interpretation.}
\begin{itemize}
    \item MSE penalizes large errors more strongly
    \item MAE is more robust to outliers
\end{itemize}

\subsection{Classification Metrics}

\medskip

\noindent
\textbf{Accuracy.}
\[
\frac{\text{number of correct predictions}}{\text{total number of predictions}}
\]

\medskip

\noindent
\textbf{Precision and recall.}
\begin{itemize}
    \item Precision: fraction of predicted positives that are correct
    \item Recall: fraction of actual positives that are detected
\end{itemize}

\medskip

\noindent
\textbf{F1 score.}
\[
\text{F1} = 2 \cdot \frac{\text{precision} \cdot \text{recall}}{\text{precision} + \text{recall}}
\]

\medskip

\noindent
\textbf{ROC and AUC.}
\begin{itemize}
    \item ROC curve plots true positive rate vs false positive rate
    \item AUC summarizes overall classification performance
\end{itemize}

\medskip

\noindent
\textbf{Imbalanced data.}  
Accuracy can be misleading when classes are imbalanced; precision, recall, and AUC provide more informative evaluation.

\medskip

\noindent
\textbf{Probabilistic metrics.}  
For models that output probabilities, cross-entropy loss is often used as an evaluation metric.

\section{Model Comparison}

Selecting a model involves comparing multiple candidates based on validation performance.

\medskip

\noindent
\textbf{Hyperparameters.}
\begin{itemize}
    \item Regularization strength
    \item Model architecture
    \item Optimization parameters
\end{itemize}

\medskip

\noindent
\textbf{Selection procedure.}
\begin{enumerate}
    \item Define a set of candidate models
    \item Evaluate each model using validation or cross-validation
    \item Select the model with best performance
\end{enumerate}

\medskip

\noindent
\textbf{Model selection vs evaluation.}
\begin{itemize}
    \item Model selection uses validation data
    \item Final evaluation uses test data
\end{itemize}

\medskip

\noindent
\textbf{Overfitting in model selection.}  
Testing many models increases the risk of selecting one that performs well by chance.

\medskip

\noindent
\textbf{Practical strategies.}
\begin{itemize}
    \item Use cross-validation
    \item Limit the search space of hyperparameters
    \item Prefer simpler models when performance is similar
\end{itemize}

\medskip

\noindent
\textbf{Tradeoffs.}
\begin{itemize}
    \item Accuracy vs complexity
    \item Performance vs computational cost
\end{itemize}

\medskip

\noindent
\textbf{Reproducibility.}  
Reliable model comparison requires consistent data splits and controlled experimental settings.

\section{A Unifying Perspective}

Model evaluation and selection are integral parts of the learning process.

\medskip

\noindent
\textbf{Pipeline.}
\begin{itemize}
    \item Train models on data
    \item Estimate generalization performance
    \item Select the best model
    \item Evaluate final performance
\end{itemize}

\medskip

\noindent
\textbf{Key challenges.}
\begin{itemize}
    \item Limited data
    \item Variance in estimates
    \item Risk of overfitting during selection
\end{itemize}

\medskip

\noindent
\textbf{Core insight.}  
Reliable evaluation requires separating the roles of training, validation, and testing.

\section{Outlook}

This chapter introduced the principles of model evaluation and selection:
\begin{itemize}
    \item separating data for unbiased evaluation,
    \item using cross-validation for robust estimates,
    \item choosing appropriate performance metrics,
    \item comparing models systematically.
\end{itemize}

\medskip

In the next chapter, we study optimization methods in greater depth, focusing on practical techniques used in large-scale machine learning.

\medskip

\noindent
\textbf{Guiding principle.}  
A model is only as good as the way it is evaluated. Reliable evaluation is essential for meaningful learning.